\chapter{Machine learning denoisers}

Also refered as data-driven denoisers, can be classified as:
\begin{enumerate}
\item Supervised (paired): supervised denoisers require one o more
  clean images and their respective noisy instances to learn, by
  comparison, what is noise and what is the clean signal (a ground
  thruth must exist). In other words, we must have one or more pairs
  $\{(\tilde{\bm{X}}_i,\bm{X}_i)\}$. Supervised denoisers are also
  classified as \gls{N2C} methods and distribution-based methods,
  because the model learns the statistical ``look'' of clean images.
\item Unsupervised (unpaired): Unsupervised denoising does not require
  labels, i.e., we don't need $(\tilde{\bm{X}},\bm{X})$ pairs,
  although we need a collection of noisy images (with the type of
  noise we want to remove) and a a collection of clean images (the
  content of the images could be any and different to the content of
  the noisy ones). So, we need a $\{\tilde{\bm{X}}_i\}_{i=0}^{N}$ and
  a $\{\bm{Y}_i\}_{i=0}^{M}$. GAN-based denoiser can be classified as
  unsupervised.
\item Self-supervised (blind / noise-only): Unsing patching of blind-spotting techniques, we simulate the existence of the pairs $\{(\tilde{\bm{X}}_i,\bm{X}_i)\}$ Learn from noisy images
  themselves. Can be:
  \begin{enumerate}
  \item Noise2Noise: Based on mapping one noisy
    version of a scene to another noisy version of the same scene.
  \item Noise2Void:
  \item Noise2Self:
  \item Blind-spot networks: predict a
    pixel from its neighbors.
  \end{enumerate}
\item ``Zero-Shot'' Denoisers: when you have only one instance
  $\tilde{\bm{X}}$ and during the training
  $model(\tilde{\bm{X}})=\tilde{\bm{X}}$, you stop when you consider
  that the model is learning the noise. For example, DIP, based on the
  fact that autoencoders with skip connections (U-Nets, although this
  is true also for all denoising \gls{DL}-based techniques) learn
  first those information in an image that is not noise, an U-Net is
  trained to map a random image $\bm{N}$ to $\tilde{\bm{X}}$. In an
  optimal point of the training process, the U-Net will generate a
  $\hat{\bm{X}\approx\bm{X}$.
\end{enumerate}


\section{Multivariate data}
%{{{

Multivariate data refers to datasets that contain more than two
variables for each observation. These variables represent different
characteristics or measurements related to the observed phenomenon. In
simpler terms, if you are collecting information about something and
recording three or more different aspects for each item, you are
dealing with multivariate data. Example: A study on students: You
might collect data on their age, test scores in math, test scores in
science, attendance rate, and extracurricular activities. Each student
is an observation, and the five pieces of information are the multiple
variables.

%}}}

\section{Training of a regression model for denoising} % Noise2Noise: Learning Image Restoration without Clean Data https://arxiv.org/pdf/1803.04189
A regression model, e.g., a \gls{CNN} can be trained using a number of pairs $(\hat{{\mathbf x}}_i, {\mathbf y}_i)$ of corrupted inputs $\hat{{\mathbf x}}_i$ and clean targets ${\mathbf y}_i$, and minimizing the empirical risk
\begin{equation}
  \underset{\theta}{\operatorname{arg\,min}} \, \sum_i L \big(f_\theta(\hat{\mathbf x}_i), {\mathbf y}_i\big)
  \label{eq:supervised_learning}
\end{equation}
where $f_\theta$ is a parametric family of mappings (e.g., CNNs) under the loss function $L$.

